\documentclass[11pt,letterpaper]{article}
\usepackage[T1]{fontenc}
\usepackage{lmodern,microtype,amsmath,amssymb,amsthm,mathtools}
\usepackage[margin=.90in,headheight=15pt]{geometry}
\usepackage{xcolor,booktabs,tabularx,array,enumitem,fancyhdr}
\usepackage[colorlinks=true,linkcolor=teal!60!black,urlcolor=teal!60!black]{hyperref}
\definecolor{ink}{HTML}{193B4C}
\definecolor{accent}{HTML}{006B70}
\pagestyle{fancy}\fancyhf{}\fancyhead[L]{\small\sffamily\color{ink}A finite-quotient proof of the bottom of the theta tower}\fancyfoot[R]{\small\thepage}
\setlength{\parindent}{0pt}\setlength{\parskip}{5pt plus 1pt minus 1pt}\setlength{\emergencystretch}{2em}
\setlist{itemsep=2pt,topsep=3pt}\raggedbottom
\newtheorem{theorem}{Theorem}[section]\newtheorem{lemma}[theorem]{Lemma}\newtheorem{proposition}[theorem]{Proposition}\theoremstyle{remark}\newtheorem{remark}[theorem]{Remark}
\newcommand{\C}{\mathbb C}\newcommand{\Ocal}{\mathcal O}\newcommand{\p}{\mathfrak p}\newcommand{\G}{\mathrm{GL}}\newcommand{\SB}{\mathcal S}\newcommand{\supp}{\operatorname{supp}}\newcommand{\Tr}{\operatorname{Tr}}\newcommand{\id}{\mathrm{id}}\newcommand{\av}{\operatorname{Av}}\newcommand{\B}{\mathcal B}
\newcommand{\baseurl}{https://github.com/VladimirReshetnikov/ProveIt/blob/9050b04000647d588d0c78d945180abcd3b0bcf7/Analysis/GeneralizedRiemannHypothesis/}
\newcommand{\src}[3]{\href{\baseurl Lean/grh/src/#1.lean\#L#2}{#3}}
\newcommand{\bottomref}[2]{\src{New/OneLift/FalseLifts/BottomOfTower}{#1}{#2}}
\newcommand{\guide}[2]{\href{\baseurl docs/guide/theta_lifts_and_dirichlet_zeros.tex\#L#1}{#2}}
\begin{document}
\thispagestyle{plain}
{\small\sffamily\color{accent}\bfseries PROOF SIMPLIFICATION \hfill 27 NOVEMBER 2026}\par\vspace{.15in}
{\LARGE\bfseries\color{ink} Eliminate the shrinking-coset argument\\ at the bottom of the theta tower}\par\vspace{.08in}
{\large A single vector-valued finite average for false-lift orthogonality}\par\vspace{.16in}
\hrule\vspace{.12in}
\textbf{Abstract.} The lower determinant estimate in the supplied Alp\"oge--Claude Lean development depends on an orthogonality theorem for false lifts. Its first local step proves vanishing against all Schwartz--Bruhat data from vanishing against a single block datum and annihilation by finite hyperplane sums. The readable implementation uses shrinking factorial-indexed cosets, a finite average of each coset indicator, a limit of local constancy, and a separate argument for cosets centered outside the good locus. We replace this passage by one finite quotient of a compact parabolic level group and a vector-valued Reynolds average. A stabilizer kills the averaged vector at every bad point. Its compact support then lies in finitely many open good orbits, where the existing block datum determines it exactly. This gives a paper proof of the same exported theorem, without changing the Hecke, Selmer, volume or analytic bounds. The proposal is distinct from the surveyed 18 proposal directories, in particular from the column--row orbital exchange and the Gram-adjoint improvement. Exact finite-field regressions and a source-level migration plan accompany the note. The Lean replacement is not compiled; the savings described below are targets, not measured reductions.

\section{Where the simplification belongs}
The overall proof compares an upper estimate and a lower estimate for the same covolume corrected by twice a local index. Its principal output is the inequality \(mR\le fa(\log b)^5\) in \src{New/Reduction/MainInequality}{75}{\texttt{New.star\_inequality}}. The upper estimate builds short integral vectors and the spanning argument supplies many independent translates. The lower estimate bounds the determinant of lifts after a suitably chosen twist. The guide describes this lower estimate at \guide{1028}{the single-lift comparison} and \guide{1204}{the role of false-lift orthogonality}. In particular, the integral comparison operator must send genuine lifts into the orthogonal complement of false lifts, so that the congruence denominator bound applies to its outputs.

The relevant route in the Lean development is
\begin{center}
\begin{tabularx}{.96\linewidth}{@{}>{\raggedright\arraybackslash}p{.38\linewidth}>{\raggedright\arraybackslash}X@{}}
\toprule
Stage & Purpose \\
\midrule
\texttt{FalseLifts/BottomOfTower} & Special block datum $\Rightarrow$ all local Schwartz--Bruhat data (this note).\\
\texttt{FalseLifts/HyperbolicPlaneStep} & Vanishing propagated up the tower using a small Hecke root.\\
\texttt{FalseLifts/PeriodsOfFalseLifts} & Vanishing periods become class sums; see \src{New/OneLift/FalseLifts/PeriodsOfFalseLifts}{127}{\texttt{descentsVanish\_hyperbolicSum}}.\\
\texttt{ComparisonOperator/GluingMaps} & Period identities make the comparison vectors orthogonal to false lifts; see \src{New/OneLift/ComparisonOperator/GluingMaps}{1056}{\texttt{isOrthogonalToFalseLifts\_glue\_periodLift}}.\\
\texttt{OneLift/LiftVolumeHeight} & Lower determinant bound using that orthogonality.\\
\bottomrule
\end{tabularx}
\end{center}
The target is \bottomref{1466}{\texttt{DescentPair.descentsVanish\_of\_hyperplaneSumsAnnihilate}}. Its proof currently passes from all data to cosets of \(i!\), then splits at every center into the integral-block-ratio case and the remaining case. The file contains 1492 physical lines. Our replacement acts directly on a general test function, instead of generating the test-function space by shrinking boxes. The arithmetic and geometric prerequisites are retained. We do not claim to remove the tower or to provide a new stand-alone RH proof.

\section{The precise local interface}
Fix a descent pair $O$ as in the source. Let $J=C\sqcup Q$ be the finite set of chosen and twist places. At each place write $F_j$ for the local field, $R_j$ for its integers, $\p_j$ for the maximal ideal, $k_j$ for its residue field. Choose the row ordering in which $O.\mathrm{incl}$ corresponds to the first $r$ rows. Use
\[
 X=\prod_{j\in J} M_{n,r}(F_j),\qquad x_j=\binom{A_j}{D_j},\qquad e=\binom{I_r}{0}.
\]
This is the finite-place portion of the existing Schwartz--Bruhat functions, which read only the components at $O.\mathrm{places}$. We retain the chosen conjugate-place conventions and the vacuum elsewhere. A point $x$ is \emph{good} if at every $j$,
\begin{equation}\label{eq:good}
 A_j\in\G_r(F_j),\qquad D_j A_j^{-1}\in M_{n-r,r}(R_j),\qquad
 D_j A_j^{-1}\in M_{n-r,r}(\p_j)\quad(j\in C).
\end{equation}
This is exactly \bottomref{266}{\texttt{HasIntegralBlockRatio}}. In particular, invertibility here is over $F_j$, not necessarily over $R_j$.

Let $H$ be the subgroup of the maximal level whose reductions have block form
\[
 \bar h_j=\begin{cases}
 \begin{pmatrix}s_j&*\\0&t_j\end{pmatrix},&j\in C,\\[3pt]
 \begin{pmatrix}s_j&0\\ *&t_j\end{pmatrix},&j\in Q.
 \end{cases}
\]
Both $s_j$ and $t_j$ are invertible. One can use the subgroup of the global maximal level defined by these conditions, acting via its local components on $X$. The local lifting lemma \bottomref{672}{\texttt{exists\_stdLevel\_atPlace\_eq}} realizes every specified tuple of local integral matrices. Let $K\le H$ be the subgroup with $s_j=I_r$ for $j\in C$. It is the original level group. Off the displayed places use exactly the source's group and vacuum conventions. In the proof below, only finite congruence quotients of $H$ occur.

Write $V=O.\mathrm{FormV}$ and $\chi_j=O.\mathrm{tamePart}(j)$. The source's \bottomref{998}{triangular-action lemma} gives the action of $H$ on $V$:
\begin{equation}\label{eq:rho}
 \rho(h)v=\left(\prod_{j\in J}\chi_j(\det s_j)^{-1}\right)
 \sigma((s_j)_{j\in C})v,
\end{equation}
where $\sigma=O.\mathrm{formAction}$ is a representation of the product of the chosen $\G_r(k_j)$'s. This is a representation because taking the diagonal block of a block-triangular matrix is multiplicative. It factors through reduction modulo the maximal ideals; on $K$ it is the scalar
\[
 \lambda(h)=\prod_{j\in Q}\chi_j(\det s_j)^{-1}.
\]
The values of $\lambda$ have absolute value one. In particular, $H$ preserves $V$ even though the whole maximal level need not do so.

For a tuple $T=(T_j)\in\prod_j\G_r(F_j)$, let $\Phi_T(x)$ denote the right translate $\Phi(xT)$ of the block datum, extended trivially at the other places as in the source. In local notation,
\begin{equation}\label{eq:phi}
 \Phi(z)=\begin{cases}
 \displaystyle\prod_{j\in Q}\chi_j(\det\bar A_j),&
 z_j\in M_{n,r}(R_j)\ \forall j,\quad \bar z_j=e\ \forall j\in C,\\
 0,&\text{otherwise},
 \end{cases}
\end{equation}
where multiplicative characters are extended by zero at zero. The right translates are Schwartz--Bruhat and are realized by source-supported finite adelic elements. Compare \bottomref{697}{\texttt{translatedBlockFn}} and \bottomref{721}{\texttt{residueBlockDatum\_mul}}.

For fixed coefficient functional $u$, archimedean polynomial $P$, and evaluation point $g$, let
\[
 B(\psi,v)=\texttt{descentPairing}\ u\ P\ g\ \psi\ v.
\]
It is linear in $\psi$ and conjugate-linear in $v$. Its equivariance, \bottomref{1103}{\texttt{descentPairing\_translate}}, is
\begin{equation}\label{eq:equiv}
 B(\psi(h^{-1}\cdot),\rho(h)v)=B(\psi,v)\quad(h\in H).
\end{equation}
The block-descent hypothesis, with all evaluation points $g$, gives $B(\Phi_T,y)=0$ for every $T$ by \bottomref{1185}{\texttt{descentPairing\_translatedBlockFn\_eq\_zero}}. The remaining hypotheses are that $\chi_j\ne 1$ for $j\in Q$, and that the source's hyperplane sums annihilate $y$. At a chosen place, using the zero row at all other places, \bottomref{1296}{\texttt{sum\_formAction\_mulSingle\_eq\_zero}} gives the single-place statement. For a nonzero row $\ell\in k_j^{1\times r}$ write
\[
 G_\ell=\{I_r+a\ell:a\in k_j^r,\ \det(I_r+a\ell)\ne0\}.
\]
Then, with all other first blocks equal to identity,
\begin{equation}\label{eq:hyper}
 \sum_{s\in G_\ell}\rho_j(s)y=0\qquad(j\in C),
\end{equation}
where $\rho_j(s)$ is the single-place action given by (\ref{eq:rho}). The source proves the twist-place counterpart on \emph{every} vector by the nontrivial determinant character. No weakening of these hypotheses is needed.

\section{A finite average with no limiting argument}
\begin{lemma}[Finite quotient for a test function]\label{lem:finite}
For $\psi\in\SB(X)$, the action of $H$ on $\psi$ and on $V$ factors, for this particular $\psi$, through a common finite quotient: more precisely, there is a normal finite-index subgroup $N\le H$ such that $N$ fixes $\psi$ and acts as the identity on $V$.
\end{lemma}
\begin{proof}
For each place choose nonnegative integers $c_j,d_j$ such that the compact support of $\psi$ lies in the product box $L=\prod_j \pi_j^{-c_j}M_{n,r}(R_j)$, and $\psi$ is invariant under additive translations by $L'=\prod_j \pi_j^{d_j}M_{n,r}(R_j)$. To see the second assertion, cover the compact support by finitely many cosets of compact open additive subgroups on which $\psi$ is constant, and take a common smaller product subgroup. If a translation moves a point outside the support into it, apply the same invariance to the reverse translation; hence invariance holds everywhere.

Take $N$ to be the kernel of the simultaneous reductions modulo $\p_j^{\max(1,c_j+d_j)}$ restricted to $H$. These rings are finite, so the quotient is finite. Each $h\in N$ preserves $L$ and satisfies $h^{-1}x-x\in L'$ on $L$. Outside $L$, both $x$ and $h^{-1}x$ are outside the support. Thus $N$ fixes $\psi$. It acts trivially on $V$ by (\ref{eq:rho}). Normality ensures that the finite quotient acts on the whole translate orbit of $\psi$.
\end{proof}
Unlike factorial-indexed cosets, these congruences are selected once per input function; no limit or inductive decomposition of small boxes is needed. One may alternatively use an existing smooth-representation finite-orbit lemma, provided it establishes the same common quotient.

\begin{lemma}[Annihilation and datum vanishing on the orbit span]\label{lem:span}
Let $U=\operatorname{span}_{\C}\{\rho(h)y:h\in H\}$. Then $U$ is finite dimensional, stable under $H$, satisfies (\ref{eq:hyper}) on every vector at every chosen place, and $B(\Phi_T,v)=0$ for all $v\in U$ and $T$.
\end{lemma}
\begin{proof}
The representation factors through a finite group, hence $U$ is the span of finitely many translates and is stable. At a chosen place the first-block action of $h$ conjugates $G_\ell$ to $G_{\ell s_j^{-1}}$, since
\[
 s_j (I+a\ell)s_j^{-1}=I+(s_j a)(\ell s_j^{-1}).
\]
The determinant weights are already incorporated in $\rho$; conjugation also preserves determinant. Consequently conjugating the annihilating sum at a nonzero row gives another sum from the same family. The action at other places commutes. Thus the sums vanish on $\rho(h)y$ and on all of $U$.

For the last assertion, fix $h\in H$. At each chosen place choose an integral invertible lift $S_j$ of its first diagonal block $s_j$; take $S_j=I$ at every twist place. Such lifts exist by lifting matrix entries (the lifted determinant is a unit), and can be made into a supported adelic element. Directly from (\ref{eq:phi}),
\begin{equation}\label{eq:right}
 \Phi_T(hx)=\left(\prod_{j\in Q}\chi_j(\det s_j)\right)\Phi_{TS}(x).
\end{equation}
Indeed, integrality is invariant under the integral invertible matrix $h_j$. At a chosen place $\bar h_j e=e s_j$, so $\bar h_j\overline{x_jT_j}=e$ is equivalent to $\overline{x_jT_jS_j}=e$. At a twist place the top block multiplies on the left by $s_j$. Applying (\ref{eq:equiv}) to $\psi=\Phi_T(h\cdot)$ gives $B(\Phi_T,\rho(h)y)=B(\Phi_T(h\cdot),y)=0$. Extend conjugate-linearly.
\end{proof}

\section{The stabilizer removes the bad locus}
\begin{lemma}[Every bad point has an annihilating stabilizer sum]\label{lem:bad}
For $x\notin X_{\mathrm{good}}$, every vector of $U$ fixed by the stabilizer $H_x=\{h\in H:hx=x\}$ is zero.
\end{lemma}
\begin{proof}
Choose a place $j$ where (\ref{eq:good}) fails. The source's \bottomref{365}{annihilating-row lemma} gives an integral row $L=(\ell_0,\ell_1)$ with $Lx_j=0$, $\bar\ell_0\ne0$, and $\bar\ell_1=0$ if $j\in Q$. For clarity, the elementary proof is as follows. If $A_j$ is singular, take a nonzero left null row of $A_j$ and append zero, dividing by a top entry of largest absolute value. If $A_j$ is invertible but a row $d$ of $D_jA_j^{-1}$ violates (\ref{eq:good}), use $L=(d,-e_i^t)$, where $e_i$ is the basis vector of that lower row, and divide by a largest top entry. Its absolute value is at least one at a chosen place and strictly larger than one at a twist place. This makes the top row primitive and, in the twist case, the bottom row zero modulo $\p_j$.

For each $s=I_r+a\bar\ell_0\in G_{\bar\ell_0}$ choose an integral lift $\tilde a$ of $a$, and define the local matrix
\[
 h_s=I_n+e\tilde a L.
\]
It fixes $x_j$ exactly. By the rank-one determinant identity its determinant reduces to $\det s\ne0$, so it is integral invertible. At a chosen place its reduction is block upper triangular; at a twist place $\bar\ell_1=0$ gives block lower triangular (in fact block diagonal). Its first diagonal block reduces to $s$. Put identity at the other places and lift to the global maximal level by the source's local placement lemma. This realizes $h_s$ as an element of $H_x$ with the single-place action $\rho_j(s)$. The construction is independent of the vector, as is visible in \bottomref{1376}{\texttt{exists\_leftMul\_eq\_and\_blockAction\_eq}}; the source's proof records the triangular reductions explicitly.

If $v\in U$ is fixed by $H_x$, then $\rho_j(s)v=v$ for every such $s$. At a chosen place, (\ref{eq:hyper}) and Lemma~\ref{lem:span} imply $|G_{\bar\ell_0}|v=0$. At a twist place, the determinant map $G_{\bar\ell_0}\to k_j^\times$ is surjective: $\det(I+a\bar\ell_0)=1+\bar\ell_0 a$ attains every nonzero value. The sum of the nontrivial character $\chi_j^{-1}\circ\det$ is zero, so the same conclusion follows (or use a single element with nontrivial character). Since the coefficient field is $\C$, $v=0$.
\end{proof}
The crucial distinction from a scalar orbit average is that the bad stabilizer acts on the value in the form space. A scalar average alone can retain the bad orbit at a chosen place; a twist by the scalar level character is not sufficient there. The enlarged group $H$, rather than only $K$, and the vector-valued average are essential.

\section{The surviving orbits are exactly block translates}
\begin{lemma}[Open good orbits]\label{lem:good}
For every good point $x$, there is a tuple $T$ with $\Phi_T(x)=1$ and $\{z:\Phi_T(z)\ne0\}=Kx$. This orbit is compact and open in $X$. For $k\in K$, $\Phi_T(kz)=\overline{\lambda(k)}\Phi_T(z)$.
\end{lemma}
\begin{proof}
Take $T_j=A_j^{-1}$. Then $x_jT_j=(I_r,D_jA_j^{-1})^t$ is integral, reduces to $e$ at chosen places, and has top determinant one at the twist places. Thus $\Phi_T(x)=1$. Any point $z$ in the nonzero support has $z_j T_j$ integral, its top block integral invertible, and at a chosen place reduction equal to $e$. Completing these columns to the square matrix
\[
 \begin{pmatrix}(z_jT_j)'&0\\ (z_jT_j)''&I_{n-r}\end{pmatrix}
\]
gives an element of $K$ by the same local placement lemma. It maps $eT_j^{-1}$ to $z_j$. The identical construction applies to $x$, hence $z\in Kx$. Conversely, covariance (\ref{eq:right}) for $k\in K$ gives $\Phi_T(kz)=\overline{\lambda(k)}\Phi_T(z)$ by $s_j=I$ at chosen places; choose $S_j=I$ there. Therefore $Kx$ is precisely the support. The support of a locally constant compactly supported function is its nonzero locus (that locus is closed as well as open), hence compact and open. These facts are also the content of \bottomref{833}{orbit-support lemma}, \bottomref{955}{normalized-datum lemma}, and \bottomref{797}{\texttt{translatedBlockFn\_leftMul}}.
\end{proof}

\section{The replacement theorem}
\begin{theorem}[Vector-valued Reynolds descent]\label{thm:main}
With the hypotheses of Section 2, for every Schwartz--Bruhat test function $\psi\in\SB(X)$, $B(\psi,y)=0$. In particular the source's exported bottom-of-tower theorem follows unchanged.
\end{theorem}
\begin{proof}
Use Lemma~\ref{lem:finite} and let $G=H/N$. For $h\in G$ write $h$ in expressions involving $\psi$ or $\rho$ using any representative. Set
\begin{equation}\label{eq:avg}
 F(x)=\frac1{|G|}\sum_{h\in G}\overline{\psi(h^{-1}x)}\,\rho(h)y\in U.
\end{equation}
This is well defined, since $N$ fixes $\psi$ and acts trivially on $V$. It is locally constant, compactly supported and equivariant: for $k\in H$, reindex $h=kh'$ in the finite sum to obtain $F(kx)=\rho(k)F(x)$. No continuity of $B$ is used.

On the space of locally constant compactly supported $U$-valued functions define the conjugate-linear contraction $\B$ by
\[
 \B\!\left(\sum_i\overline{\psi_i(\cdot)}v_i\right)=\sum_i B(\psi_i,v_i).
\]
This is well defined: choose a basis $v_i$ of the finite-dimensional space $U$ and use the unique scalar coordinate functions. Under any finite expansion, linearity in the scalar-function argument and conjugate-linearity in the vector argument give the displayed formula; equivalently it is the map on the conjugate of $\SB(X)\otimes \overline U$ using the pairing. In particular, by (\ref{eq:equiv}) and the real normalization factor,
\begin{equation}\label{eq:contract}
 \B(F)=\frac1{|G|}\sum_{h\in G} B(\psi(h^{-1}\cdot),\rho(h)y)=B(\psi,y).
\end{equation}
If $x$ is bad, equivariance makes $F(x)$ fixed by $H_x$, so Lemma~\ref{lem:bad} shows that $F(x)=0$. The nonzero support of $F$ is compact and lies in the good locus. It is $K$-invariant since $\rho(k)$ is invertible. By Lemma~\ref{lem:good}, the good $K$-orbits form a disjoint open cover of this compact support. Only finitely many, say $Kx_1,\ldots,Kx_t$, are needed. For each choose $T_i$ with $\Phi_{T_i}(x_i)=1$. At $kx_i$ we have
\[
 F(kx_i)=\lambda(k)F(x_i)=\overline{\Phi_{T_i}(kx_i)}F(x_i).
\]
Therefore, exactly as functions on all of $X$,
\begin{equation}\label{eq:decomp}
 F=\sum_{i=1}^t\overline{\Phi_{T_i}}\,F(x_i).
\end{equation}
By Lemma~\ref{lem:span}, each summand has zero contraction. Combining (\ref{eq:decomp}) with (\ref{eq:contract}) proves the claim. The case of empty support is the empty sum.

For the Lean interface, apply this for each $u,P,g$ after using \texttt{DescentsVanish.of\_schrodingerOp} exactly as at the beginning of the current theorem. The vanishing for all $g$ is extensional equality of the descent function with zero. The outer reduction to polynomial Fock vectors and vacuum components, and all subsequent arguments, are unchanged.
\end{proof}

\paragraph{The new mechanism in one equation.} Let $\pi_{\rm bad}$ denote restriction to the bad locus. The proof establishes $\pi_{\rm bad}\av_H(\bar\psi\,y)=0$, not just $B(\psi\,1_{\rm bad},y)=0$ (the latter restriction need not itself be Schwartz--Bruhat). It never needs the bad locus to be open. This is important at rank-deficient points, which can be limits of open orbits. Compactness is used \emph{after} killing them, and only then to obtain finitely many open orbits.

\section{Why this should shorten the readable Lean development}
The old and new proofs have the same local arithmetic core: an annihilating integral row gives a rank-one shear, and on a good chart completing the columns proves transitivity. In the replacement, these two pointwise facts feed into one algebraic average. The old general-purpose infrastructure for factorial cosets is no longer needed on this route. The following ranges are physical line ranges in blob \texttt{b113ad2b...} (not counts of elaborated proof terms).
\begin{center}
\begin{tabularx}{\linewidth}{@{}p{.18\linewidth}p{.13\linewidth}X@{}}
\toprule
Old range & Gross lines & Proposed change\\
\midrule
387--557 & 171 & Nested cosets, finite means, factorial boxes, eventual constancy and coset-span lemma; replace by uniform congruence quotient.\\
624--645 & 22 & Translation of coset indicators, including fixed-center case; candidate removal if not needed elsewhere.\\
886--954 & 69 & Scalar level translation/finite orbit and proportionality of averaged coset; replace by $H$-equivariance and open-orbit decomposition.\\
1211--1263 & 53 & Idempotent-component bookkeeping and separate good-coset argument; replace by Lemma~\ref{lem:good} and (\ref{eq:decomp}).\\
1437--1485 & 49 & Separate bad-coset theorem and final eventual-coset glue; replace by stabilizer lemma and Theorem~\ref{thm:main}.\\
\midrule
Total & 364 & Gross affected region, before adding new code or retaining shared declarations.\\
\bottomrule
\end{tabularx}
\end{center}
A code-search check found the symbols \texttt{mem\_of\_nested\_cosets}, \texttt{factorialBox}, \texttt{factorialIdele}, \texttt{mem\_of\_eventually\_coset}, \texttt{finite\_range\_levelTranslate}, \texttt{eq\_smul\_translatedBlockFn} and the two separate coset-vanishing lemmas only in this file. Search is not a compiler dependency analysis; other helpers may be shared. In particular retain \texttt{HyperplaneSumsAnnihilate}, its constructor-facing API, the scalar group action identity if it is useful elsewhere, and all the local shear and block-completion algebra. The 364 lines are about 24\% of this particular file, not of the entire project. Net savings cannot be measured until the new generic finite quotient and vector-valued contraction are formalized. The payoff is stronger than a cosmetic tactic rewrite: the exported theorem no longer needs to reason about an arbitrary shrinking sequence of neighborhoods. A reusable finite-orbit lemma can be placed in a library module.

\paragraph{Migration.} (1) Introduce the parabolic subgroup and its representation using \texttt{blockAction} and \texttt{smul\_chosenGLRep\_eq\_apply\_of\_triangular}; record the single-place first-block projection and its conjugation law. (2) Establish a single simultaneous congruence quotient fixing the given Schwartz--Bruhat function. Reuse existing finite-place support and local constancy API rather than recreating Haar measure. (3) Define the finite-dimensional orbit span and use the conjugate-linear pairing through basis coordinates, or an equivalent tensor-product construction. (4) Generalize the shear witness to include its triangular membership; no new local estimates. (5) Prove (\ref{eq:right}) from the block datum. (6) Use the support-orbit lemma and compactness to write (\ref{eq:decomp}), then replace the old theorem proof. See the accompanying \texttt{LEAN\_PLAN.md} for implementation details and acceptance tests.

\section{Nonduplication audit and limitations}
The three supplied directories were accessed using the GitHub connector. The target, guide and periods file were confirmed at revision \texttt{9050b04000647d588d0c78d945180abcd3b0bcf7}. The proposal directory had 18 subdirectories. Their recursive listings and principal TeX sections or abstracts were examined; an earlier detailed five-directory audit was also read. Targeted scans of the main sources for bottom-of-tower names, hyperplane sums, false lifts and coinvariants found no prior proposal replacing this passage. This is an overlap assessment against the inspected sources, not an exhaustive literature priority search or a complete line-by-line audit of every attachment. See \texttt{proposal\_audit.md} for the directory-by-directory matrix.

The closest existing proposal, \texttt{orbital\_theta\_exchange}, proves a scalar twisted-coinvariant theorem for matching column and row data in the determinant exchange. Its reusable finite-orbit averaging ideas may be assumed implemented. The present theorem targets a different interface: a stabilizer acts on an orbit of \emph{form vectors}, the annihilation condition for nonzero hyperplane rows excludes a locus, and the good support is covered by finitely many open orbits. It requires neither complementary frames nor a column--row partner. The Gram-adjoint proposal explicitly retains the false-lift orthogonality theorem; our replacement reduces its proof burden rather than recomputing the comparison operator. The earlier finite-label proposals may change how the hyperplane hypothesis is established, but the statement here takes exactly that hypothesis as input. Thus these changes are composable.

The mathematical derivation above is a proof of a local replacement conditional on the source's stated interfaces. It is not a new independent certification of the full RH/GRH development. No Lean executable was available in the artifact environment; no claim of compilation, axiom audit or actual line reduction is made. No changes were made to the repository. Before accepting the patch, the original theorem statement and its downstream consumers must build with the original toolchain and axioms. Do not introduce any new axiom, \texttt{sorry}, or opaque proof assumption to supply the averaging interface.

\section{Verification and further research}
The included \texttt{check\_average.py} checks exact finite-field models of the stabilizer and averaging statements. It exhausts all delta functions on rank-one matrices over $\mathbb F_3$ and rank-two matrices over $\mathbb F_2$ (with $n=r+1$). It verifies equivariance, vanishing on all bad points, and the orbit decomposition geometry; the chosen rank-two example uses the sign of the permutation on the three nonzero vectors of $\mathbb F_2^2$. The scalar-character examples also deliberately omit annihilation or use a trivial twist character; the bad support then survives. These are finite regressions, not tests of the adelic theorem or of a general higher-dimensional representation. The proof of Theorem~\ref{thm:main}, rather than the regressions, covers arbitrary rank and local fields.

Three follow-on questions are worth testing after the patch. First, can the entire bottom step be packaged as a generic equivariant compact-support theorem over a profinite group with only open good orbits and vanishing stabilizer invariants, then reused in other theta descent proofs? Second, can the hyperplane condition be supplied directly in the quotient representation by a cuspidality certificate, without constructing each group-algebra element individually? This must be compared with the already planned robust-label proposals. Third, can a similar support-annihilation theorem replace parts of the higher hyperbolic-plane step, rather than merely reusing its analytic estimates? These are research directions, not results proved in this note.

\end{document}
